% !TEX root = ../main.tex
\corrigesChapitre{6}{Trigonométrie}
\corrigesSouples


\begin{blocLivre}{cor-6-1}
\begin{corrige}{1}{Convertir des angles}
\demarche{Un demi-tour vaut 180 degrés, soit pi radians.}
\begin{questions}
\item Pour convertir des degrés en radians, on multiplie par $\pi/180$ :
$30\pi/180=\res{\pi/6}$ ; $45\pi/180=\res{\pi/4}$ ;
$120\pi/180=\res{2\pi/3}$ ; $270\pi/180=\res{3\pi/2}$.
\item Dans l’autre sens, on multiplie par $180/\pi$ :
$(\pi/3)(180/\pi)=\res{60^\circ}$ ;
$(5\pi/6)(180/\pi)=\res{150^\circ}$ ;
$(7\pi/4)(180/\pi)=\res{315^\circ}$.
\end{questions}

\end{corrige}
\end{blocLivre}
\begin{blocLivre}{cor-6-2}
\begin{corrige}{2}{Même point-image}
\demarche{Deux mesures donnent le même point lorsque leur différence est un nombre entier de tours.}
\begin{questions}
\item La différence vaut $4\pi=2\times2\pi$ : oui.
\item $3\pi/2-(-\pi/2)=2\pi$ : oui. En revanche $5\pi/4-\pi/4=\pi$ : les points sont opposés, donc distincts.
\end{questions}
\label{figlivre-r06-02-images-cor}
\figureLivre[corrige]{r06-02-images}

\end{corrige}
\end{blocLivre}
\begin{blocLivre}{cor-6-3}
\begin{corrige}{3}{Lire les valeurs remarquables}
\demarche{Sur le cercle unité, l’abscisse est le cosinus et l’ordonnée est le sinus.}
\begin{questions}
\item Les coordonnées sont $\res{(1;0)}$, $\res{(\sqrt3/2;1/2)}$, $\res{(\sqrt2/2;\sqrt2/2)}$ et $\res{(0;1)}$.
\item On lit les points sur les axes : $\res{-1;\ 0;\ 0;\ -1}$ respectivement.
\end{questions}
\label{figlivre-r06-03-valeurs-cor}
\figureLivre[corrige]{r06-03-valeurs}

\end{corrige}
\end{blocLivre}
\begin{blocLivre}{cor-6-4}
\begin{corrige}{4}{Utiliser une symétrie}
\demarche{Les symétries conservent une coordonnée et changent le signe de l’autre.}
\begin{questions}
\item La réflexion par rapport à l’axe horizontal conserve l’abscisse et oppose l’ordonnée.
Ainsi $\cos(-\pi/3)=\cos(\pi/3)=\res{1/2}$ et
$\sin(-\pi/3)=-\sin(\pi/3)=\res{-\sqrt3/2}$.
\item Comme $2\pi/3=\pi-\pi/3$, la réflexion par rapport à l’axe vertical oppose l’abscisse et conserve l’ordonnée.
Donc $\cos(2\pi/3)=\res{-1/2}$ et $\sin(2\pi/3)=\res{\sqrt3/2}$.
Le point du deuxième quadrant a bien un cosinus négatif et un sinus positif.
\end{questions}

\label{figlivre-r06-04-symetries-cor}
\figureLivre[corrige]{r06-04-symetries}

\end{corrige}
\end{blocLivre}
\begin{blocLivre}{cor-6-5}
\begin{corrige}{5}{Une longueur d'arc}
\demarche{Utiliser les radians dans la formule longueur égale rayon multiplié par angle.}
La formule $L=R\theta$ donne $\res{2\pi\text{ cm}}$ et $\res{9\pi/2\text{ cm}}$. Pour un tour entier, $L=6\times2\pi=\res{12\pi\text{ cm}}$.
\label{figlivre-r06-05-arcs-cor}
\figureLivre[corrige]{r06-05-arcs}

\end{corrige}
\end{blocLivre}
\begin{blocLivre}{cor-6-6}
\begin{corrige}{6}{Réduire des grandes mesures}
\demarche{Ajouter ou retirer des tours complets avant de lire une valeur remarquable.}
\begin{questions}
\item On choisit les nombres entiers de tours qui ramènent les mesures dans $\mathopen{]}-\pi;\pi]$ :
\[\frac{17\pi}6-2\pi=\frac{5\pi}6,\qquad
-\frac{11\pi}4+2\pi=-\frac{3\pi}4,\qquad
\frac{29\pi}3-10\pi=-\frac\pi3.\]
\item Les points-images sont conservés. Les symétries et les valeurs remarquables donnent, dans le même ordre,
\[\res{\left(-\frac{\sqrt3}2;\frac12\right)},\quad
\res{\left(-\frac{\sqrt2}2;-\frac{\sqrt2}2\right)},\quad
\res{\left(\frac12;-\frac{\sqrt3}2\right)}.\]
Chaque couple est écrit dans l’ordre $(\cos;\sin)$ ; les signes se contrôlent avec le quadrant.
\end{questions}

\end{corrige}
\end{blocLivre}
\begin{blocLivre}{cor-6-7}
\begin{corrige}{7}{Retrouver une coordonnée}
\demarche{Calculer le carré du sinus, puis choisir son signe avec le quadrant.}
\begin{questions}
\item $\sin^2t=1-9/25=16/25$. Dans le deuxième quadrant, le sinus est positif : $\res{\sin t=4/5}$.
\item Pour $-t$ : $\res{(-3/5;-4/5)}$ ; pour $t+\pi$ : $\res{(3/5;-4/5)}$.
\end{questions}
\label{figlivre-r06-07-coordonnee-cor}
\figureLivre[corrige]{r06-07-coordonnee}

\end{corrige}
\end{blocLivre}
\begin{blocLivre}{cor-6-8}
\begin{corrige}{8}{Deux angles associés}
\demarche{Repérer le quadrant pour contrôler les signes des deux coordonnées.}
\begin{questions}
\item Les deux angles sont dans le troisième quadrant. Les valeurs sont $-\sqrt2/2$, $-\sqrt2/2$, $-\sqrt3/2$ et $-1/2$.
\item $A=-\sqrt2-\sqrt2/2=\res{-3\sqrt2/2}$.
\end{questions}
\end{corrige}
\end{blocLivre}
\begin{blocLivre}{cor-6-9}
\begin{corrige}{9}{Retrouver des angles sur le cercle}
\demarche{Une condition sur le cosinus fixe une abscisse ; une condition sur le sinus fixe une ordonnée.}
\begin{questions}
\item Les abscisses égales à $1/2$ correspondent à $\res{t=\pi/3\text{ ou }5\pi/3}$.
\item Les ordonnées égales à $-\sqrt2/2$ donnent $\res{t=5\pi/4\text{ ou }7\pi/4}$.
\item Les points sur l’axe vertical donnent $\res{t=\pi/2\text{ ou }3\pi/2}$.
\end{questions}

\label{figgeo-06-09-cor}
\figureGeo[corrige]{06-09}
\end{corrige}
\end{blocLivre}
\begin{blocLivre}{cor-6-10}
\begin{corrige}{10}{Une portion de cercle}
\demarche{Calculer la fraction d’un tour avant l’aire et le périmètre du secteur.}
\begin{questions}
\item La fraction est $(2\pi/5)/(2\pi)=\res{1/5}$.
\item L’aire vaut $100\pi/5=\res{20\pi\text{ cm}^2}$. L’arc mesure $10\times2\pi/5=4\pi$ cm ; le périmètre est $\res{20+4\pi\text{ cm}}$.
\end{questions}

\label{figgeo-06-10-cor}
\figureGeo[corrige]{06-10}
\end{corrige}
\end{blocLivre}
\begin{blocLivre}{cor-6-11}
\begin{corrige}{11}{Une identité à démontrer}
\demarche{Développer les deux carrés et regrouper les termes opposés.}
\begin{questions}
\item Le premier carré vaut $\cos^2t+2\cos t\sin t+\sin^2t$ ; le second vaut
$\cos^2t-2\cos t\sin t+\sin^2t$. En les additionnant, les deux termes croisés se compensent.
\item La somme devient $2\cos^2t+2\sin^2t=2(\cos^2t+\sin^2t)=\res{2}$,
car le point du cercle unité vérifie $\cos^2t+\sin^2t=1$.
\item Pour $t=\pi/4$, le sinus et le cosinus valent $\sqrt2/2$.
Le premier carré vaut $(\sqrt2)^2=2$ et le second vaut $0^2=0$, ce qui confirme le résultat.
\end{questions}

\end{corrige}
\end{blocLivre}
\begin{blocLivre}{cor-6-12}
\begin{corrige}{12}{Une grande roue}
\demarche{Ajouter l’altitude du centre à la composante verticale du rayon.}
\begin{questions}
\item On obtient $\res{(12;14)}$, $\res{(0;26)}$, $\res{(-12;14)}$ et $\res{(0;2)}$.
\item Les hauteurs sont $14+12/2=\res{20\text{ m}}$ et $14-12/2=\res{8\text{ m}}$.
\item La durée est $48\times(5\pi/6)/(2\pi)=\res{20\text{ s}}$, sous l’hypothèse d’une rotation uniforme.
\end{questions}

\label{figgeo-06-12-cor}
\figureGeo[corrige]{06-12}
\end{corrige}
\end{blocLivre}
\begin{blocLivre}{cor-6-13}
\begin{corrige}{13}{Un hexagone régulier}
\demarche{Découper l’hexagone en six triangles centraux équilatéraux.}
\begin{questions}
\item Les angles sont $k\pi/3$ pour $k=0,\ldots,5$. Les coordonnées sont $(1;0)$, $(1/2;\sqrt3/2)$, $(-1/2;\sqrt3/2)$, $(-1;0)$, $(-1/2;-\sqrt3/2)$, $(1/2;-\sqrt3/2)$.
\item Le carré de la distance vaut $(1/2)^2+(\sqrt3/2)^2=1$ : chaque côté mesure $\res{1}$.
\item Le périmètre vaut $\res{6}$. Les six triangles équilatéraux ont chacun pour aire $\sqrt3/4$ : aire totale $\res{3\sqrt3/2}$.
\end{questions}

\label{figgeo-06-13-cor}
\figureGeo[corrige]{06-13}

Pour contrôler l’aire, chaque triangle équilatéral de côté 1 a pour hauteur $\sqrt{1-(1/2)^2}=\sqrt3/2$. Son aire est donc $\sqrt3/4$, à multiplier par six.

\end{corrige}
\end{blocLivre}
\begin{blocLivre}{cor-6-14}
\begin{corrige}{14}{Une zone autorisée}
\demarche{Traduire les deux contraintes par des zones du cercle, puis prendre leur intersection.}
\begin{questions}
\item Le point est dans le demi-plan supérieur : $t\in[0;\pi]$.
\item Le point est à gauche de la droite $x=-1/2$ : $t\in[2\pi/3;4\pi/3]$.
\item L’intersection est $\res{[2\pi/3;\pi]}$. L’angle parcouru vaut $\pi/3$ ; sur le cercle unité, la longueur de cet arc est donc $\res{\pi/3}$.
\end{questions}

\label{figgeo-06-14-cor}
\figureGeo[corrige]{06-14}
\end{corrige}
\end{blocLivre}
\begin{blocLivre}{cor-6-15}
\begin{corrige}{15}{Approcher un cercle par des polygones}
\demarche{Partager un triangle central en deux triangles rectangles.}
\begin{questions}
\item Le rayon du cercle unité vaut 1. Dans le demi-triangle rectangle, le demi-côté vaut
$\sin(\pi/n)$ : un côté entier mesure $2\sin(\pi/n)$. Pour $n$ côtés,
$\res{P_n=2n\sin(\pi/n)}$.
\item $P_6=12\sin(\pi/6)=\res{6}$.
Avec la valeur fournie, $P_{12}=24(\sqrt6-\sqrt2)/4=\res{6(\sqrt6-\sqrt2)}\approx6{,}21166$.
\item Le cercle a pour longueur $2\pi\approx6{,}28319$.
Les deux périmètres calculés sont plus petits ; celui à douze côtés est plus proche.
Ces deux comparaisons numériques ne démontrent pas une convergence pour tous les $n$.
\item La relation admise double le nombre de côtés sans utiliser $\pi$ :
\end{questions}
\begin{python}
from math import sqrt
def archimede(k):
    n, c = 6, 1.0
    for i in range(k):
        c = sqrt(2 - sqrt(4 - c*c))
        n = 2*n
    return n*c/2
\end{python}
Pour $k=1$, $c_{12}=\sqrt{2-\sqrt3}=(\sqrt6-\sqrt2)/2$ : on retrouve $P_{12}$.
Pour $k=4$, le polygone a 96 côtés et la fonction renvoie environ $\res{3{,}14103}$.
Chaque côté est plus court que l'arc qu'il sous-tend, donc $P_n<2\pi$ : ce demi-périmètre
est une valeur approchée par défaut de $\pi$ ;
Archimède en déduisait $\pi>3+10/71$. La relation admise se démontre avec Al-Kashi :
$c_{2n}^2=2-2\cos(\pi/n)$ et $\cos(\pi/n)=\sqrt{1-c_n^2/4}$.
\label{figgeo-06-15-cor}
\figureGeo[corrige]{06-15}
\end{corrige}
\end{blocLivre}
